% !TEX root = ../main.tex
%% Corrigé du devoir surveillé n° 3 — valeurs vérifiées : verifications/devoirs.py

\corrigeDevoir{3}

\begin{corrigeBac}{8}{L'arôme de banane}
\begin{questions}
\item Acide éthanoïque $\ce{CH3-COOH}$ : groupe \res{\text{carboxyle}}, famille des \res{\text{acides carboxyliques}}.\newline
3-méthylbutan-1-ol $\ce{(CH3)2CH-CH2-CH2-OH}$ : groupe \res{\text{hydroxyle}}, famille des \res{\text{alcools}}.\newline
Ester $\ce{CH3-COO-CH2-CH2-CH(CH3)2}$ : groupe \res{\text{ester}} $\ce{-COO-}$, famille des \res{\text{esters}}. \pts{1,5 pt}
\item Ballon chauffé par un chauffe-ballon, surmonté d'un réfrigérant à eau (schéma ci-dessous). Le chauffage \res{\text{accélère}} la
transformation ; le réfrigérant condense les vapeurs, qui retombent dans le ballon : \res{\text{aucune perte de matière}}. \pts{1,5 pt}
\item $n(\text{alcool})=\dfrac{\rho\,V}{M}=\dfrac{0{,}81\times21{,}7}{88{,}0}=\res{\qty{0,200}{mol}}$ ; $n(\text{acide})=\dfrac{1{,}05\times23{,}0}{60{,}0}=\res{\qty{0,40}{mol}}$. Les
réactifs réagissent mole à mole : l'\res{\text{alcool}} est le réactif limitant. \pts{1,5 pt}
\item $n_\mathrm{max}(\text{ester})=\qty{0,200}{mol}$, soit $m_\mathrm{max}=0{,}200\times130{,}0=\res{\qty{26,0}{g}}$ ; rendement : $\eta=\dfrac{20{,}8}{26{,}0}=\res{\qty{80}{\percent}}$. \pts{1,5 pt}
\item Le spectre \res{\text{B}} : il présente une bande forte et fine vers \qty{1740}{cm^{-1}} (C$=$O d'un ester) et aucune bande O--H. Le
spectre A, avec sa large bande vers \qty{3340}{cm^{-1}} (O--H lié), est celui de l'alcool. \pts{1,5 pt}
\item La transformation n'est pas totale : l'excès d'acide, réactif bon marché et facile à éliminer par lavage, permet de
\res{\text{consommer davantage d'alcool}} et d'améliorer le rendement. \pts{0,5 pt}
\end{questions}
\schema[corrige]{09-reflux}
\schema[corrige]{ds3-ir}
\end{corrigeBac}

\begin{corrigeBac}{7}{Le réchaud de camping}
\begin{questions}
\item $\res{\ce{2 C4H10 + 13 O2 -> 8 CO2 + 10 H2O}}$, soit, pour une mole de butane : $\ce{C4H10 + 13/2 O2 -> 4 CO2 + 5 H2O}$. \pts{1 pt}
\item Liaisons rompues : $3\,\text{C--C}+10\,\text{C--H}+6{,}5\,\text{O=O}$, soit $3\times345+10\times415+6{,}5\times498=\qty{8422}{kJ}$ reçus. Liaisons formées :
$8\,\text{C=O}+10\,\text{O--H}$, soit $8\times804+10\times463=\qty{11062}{kJ}$ cédés. $E_\mathrm{r}=8422-11062=\res{\qty{-2640}{kJ.mol^{-1}}}$ : la
combustion libère \qty{2640}{kJ} par mole de butane. \pts{2 pts}
\item $\dfrac{\qty{2640}{kJ}}{\qty{0,0580}{kg}}=\qty{4,55e4}{kJ.kg^{-1}}=\res{\qty{45,5}{MJ.kg^{-1}}}$ : l'estimation est \res{\text{en accord}} avec la valeur des
tables (\qty{45,8}{MJ.kg^{-1}}, soit moins de \qty{1}{\percent} d'écart). \pts{1 pt}
\item $E=0{,}190\times45{,}8=\res{\qty{8,70}{MJ}}$. \pts{1 pt}
\item Énergie utile : $0{,}40\times8{,}70=\qty{3,48}{MJ}$. Pour \qty{1,0}{L} d'eau : $4{,}18\times1{,}0\times(100-20)=\qty{334}{kJ}$. On peut donc chauffer
$\dfrac{3{,}48\times10^{3}}{334}\approx\res{\qty{10}{L}}$ d'eau (10,4 L). \pts{1,5 pt}
\item $n(\ce{C4H10})=\dfrac{190}{58{,}0}=\qty{3,28}{mol}$, donc $n(\ce{CO2})=4\times3{,}28=\qty{13,1}{mol}$ et $m=13{,}1\times44{,}0\approx\res{\qty{577}{g}}$. \pts{0,5 pt}
\end{questions}
\retiens{Énergie de réaction = énergie des liaisons rompues -- énergie des liaisons formées : négative pour une combustion, qui libère de l'énergie.}
\end{corrigeBac}

\begin{corrigeBac}{5}{Le bioéthanol}
\begin{questions}
\item $\ce{CH3-CH2-OH}$ : groupe \res{\text{hydroxyle}}, famille des \res{\text{alcools}}. \pts{1 pt}
\item L'éthanol présente une \res{\text{bande large}} vers 3200 -- \qty{3400}{cm^{-1}} (O--H lié) ; le méthoxyméthane, sans liaison O--H,
\res{\text{n'en a pas}}. \pts{1 pt}
\item $\res{\ce{C2H6O + 3 O2 -> 2 CO2 + 3 H2O}}$. \pts{1 pt}
\item Éthanol : \qty{1}{MJ} demande $\dfrac{1000}{26{,}8}=\qty{37,3}{g}$, soit \qty{0,811}{mol}, qui rejette $2\times0{,}811$ mol de $\ce{CO2}$ : \res{\qty{71}{g}}. Butane :
$\dfrac{1000}{45{,}8}=\qty{21,8}{g}$, soit \qty{0,376}{mol}, qui rejette $4\times0{,}376$ mol de $\ce{CO2}$ : \res{\qty{66}{g}}. L'éthanol rejette
\res{\text{un peu plus}} de dioxyde de carbone par mégajoule. \pts{1,5 pt}
\item Pourtant, le carbone rejeté par le bioéthanol a été \res{\text{prélevé dans l'atmosphère}} par les plantes lors de leur croissance
(photosynthèse) ; celui du butane, d'origine fossile, s'ajoute au $\ce{CO2}$ de l'atmosphère. \pts{0,5 pt}
\end{questions}
\end{corrigeBac}
